% ======================================================================
% LAW CHANGE NOTE, 2026-08-03. See ../docs/LAW_CHANGE.md.
%
% THE DENOMINATOR DEFINED BELOW IS SUPERSEDED, and this is a separate change
% from the illumination law itself.
%
% chi is defined here as the arriving power density over what THE SAME NETWORK
% would deliver at the same point OVER OPEN GROUND. That needs the network to
% still exist once the buildings are removed. Under the replacement law the
% sites sit ON the buildings, on the facade tips, so removing the buildings
% removes the network and the denominator is empty.
%
% The replacement is to report the geometric quantity per unit site density and
% per unit transmit power, with a second reference taken against the unoccluded
% roofline.
%
% Two things that DO survive:
%  - The argument for reporting a ratio at all. The absolute level is still
%    unknowable and a common multiplier still cancels. Site density and
%    transmit power enter as one multiplier in front of every city's answer,
%    so neither ever has to be stated or defended.
%  - The shadowing result, which is the most robust thing in the study because
%    it needs no photographs and holds at all eleven cities.
%
% Related: normalising the line of sight term to one is a REPORTING SPLIT, not
% a definition of source power. Reading it as a definition would mean operators
% radiate harder where there is less sky, and it would delete the shadowing
% result. The answer is reported as two factors, the line of sight term and the
% multipath surplus on top of it.
%
% Nothing below has been edited.
% ======================================================================

\section{Exposure ratio}
\label{sec:ratio}

A person standing in a city square receives radio power from many base stations
at once. The buildings around them block part of that power and reflect part of
it back, so the total depends on the shape of the square as much as on any
single antenna. This paper computes that dependence as one number,
\begin{equation}
  \chi = \frac{S}{S_0} ,
  \label{eq:chidef}
\end{equation}
where $S$ is the power density arriving at the head in the real square and $S_0$
is the power density the same network would deliver at the same point over open
ground. This ratio is called the exposure ratio.

The ratio is the useful quantity because the network is unknown. An absolute
exposure in W/m$^2$ needs the transmit power, the number of sites and their
positions, none of which is public at millimeter wave frequencies and none of
which is a property of the square. Dividing by the free space answer removes all
three. What is left is a pure number that is 1 when the square is empty, below 1
where the buildings shadow more than they reflect, and above 1 where they
reflect more than they shadow. The ratio transfers between cities, which a power
density does not.

Computing $\chi$ directly would mean placing a transmitter at every position a
base station could occupy and tracing rays from each one to the person, which is
the expensive direction to work in. Instead everything is computed at the one
point where the person stands. Rays are launched from that point, followed until they leave
the square, and weighted by how likely a base station is to sit in the direction
each one left. \Cref{fig:adjoint} shows the idea.

Putting the transmitter and the receiver at the same point is what makes the
rest of the method possible, because a street level photograph taken from that
same point sees the surfaces the early bounces will hit. The material of a
facade is therefore measured rather than assumed, and how far that measurement
reaches is what sets how many bounces are worth following.
Section~\ref{sec:budget} returns to this and fixes the bounce budget from it.

\begin{figure*}[!t]
  \centering
  \includegraphics[width=\textwidth]{20_adjoint_idea}
  \caption{Rays are shot outward from the person rather than inward from the
  sources. A path is reversible, so a ray that leaves in direction $\uloc$ and
  escapes the square in direction $\uext$ carries the same attenuation as a ray
  arriving from a source in direction $\uext$ and reaching the person from
  direction $\uloc$.}
  \label{fig:adjoint}
\end{figure*}

\subsection{Notation}

Vectors are boldface, unit vectors carry a hat. $\uloc$ is a direction leaving
the observation point and $\uext$ is the direction in which a ray finally
escapes. Angles are measured from the horizon upward, so $\alpha$ is elevation
and $\alpha = 0$ is level with the head. \Cref{tab:symbols} lists the symbols
used in the equations.

\begin{table}[!t]
  \caption{Symbols}
  \label{tab:symbols}
  \centering
  \begin{tabular}{llc}
    \toprule
    Symbol & Meaning & Unit \\
    \midrule
    $\x$ & observation point, the head of the pedestrian & m \\
    $\uloc$ & direction leaving $\x$ & -- \\
    $\uext$ & direction in which a ray escapes the scene & -- \\
    $\alpha$ & elevation above the horizon & rad \\
    $r$ & horizontal range from $\x$ to a site & m \\
    $\rho$ & slant range from $\x$ to a site & m \\
    $h$ & height of a site above $\x$ & m \\
    $\nu$ & number of sites per unit volume & m$^{-3}$ \\
    $Q(\uext)$ & share of incident power per steradian of sky & sr$^{-1}$ \\
    $\Phi(\uloc)$ & power arriving per steradian, in units of $S_0$ & sr$^{-1}$ \\
    $S$, $S_0$ & arriving and free space power density & W/m$^2$ \\
    $\chi$ & exposure ratio, $1$ in free space & -- \\
    $w$ & power carried by one ray, $1$ at launch & -- \\
    $\Gamma$ & Fresnel reflection coefficient & -- \\
    $\theta_i$ & angle of incidence at a surface & rad \\
    $\kappa$ & fraction of reflected power that stays specular & -- \\
    $\sigma_h$ & root mean square surface roughness & m \\
    $\varepsilon_r$ & complex relative permittivity & -- \\
    $\lambda$ & wavelength & m \\
    $N$ & rays launched per standpoint & -- \\
    $M$ & cells in the launch grid & -- \\
    $L$ & largest number of surface interactions & -- \\
    $S_{\mathrm{inc}}$ & power density incident on a body facet & W/m$^2$ \\
    $S_{ab}$ & absorbed power density on the body & W/m$^2$ \\
    $T_0$ & power transmitted into tissue at normal incidence & -- \\
    $\nhat$ & outward normal of a body facet & -- \\
    $\khat$ & propagation direction of an incident wave & -- \\
    \bottomrule
  \end{tabular}
\end{table}

\subsection{Assumptions}

The method rests on seven assumptions, stated here because each one bounds what
the results can mean.

\begin{enumerate}
  \item[A1] \emph{Far field.} Every source is far enough away that it
  illuminates the square with a plane wave. The closest source position any
  model admits is 10\,m from the head, at the near edge of the street level
  band, and the array apertures involved are under 0.3\,m, whose far field
  begins at $2D^2/\lambda = 9$\,m at 15\,GHz. The margin is thin only for the
  nearest street level sources and is wide everywhere else.
  \item[A2] \emph{Powers add, phases do not.} Contributions from different
  paths are summed in power. At 15\,GHz and a 100\,m range, the phase of a path
  turns by about $3.1 \times 10^{4}$\,rad per radian of change in the exit
  angle, so resolving it would need on the order of $10^{10}$ angular cells
  against the 512 used here. Any phase this method could carry would be noise.
  \item[A3] \emph{Sites are a density, not a list.} Base station positions are
  unknown and are modeled as a uniform density over a slab of ground, bounded
  in height and in range. Section~\ref{sec:illum} builds the model.
  \item[A4] \emph{Every site delivers the same power at the head.} Sites are
  counted rather than weighted by their range, so the angular distribution of
  incident power follows the angular distribution of site count. The reason is
  that each site is built to serve its own cell and sets its power to cover it,
  so a nearer site is not a proportionally louder one. That is an assumption
  about how a network is operated rather than about how waves travel, and
  Section~\ref{sec:illum} carries the alternative through as a sensitivity. The
  choice is not cosmetic: the street level
  model places 88\% of its measure below $5^\circ$ under counting and 42\%
  under spreading, and elevation weighting is what separates one square from
  another.
  \item[A5] \emph{Surfaces are locally flat.} Each triangle is treated as a
  smooth interface with a Gaussian height error of standard deviation
  $\sigma_h$ about its own plane.
  \item[A6] \emph{No diffraction.} Energy bending around building edges is
  absent. The paragraphs after this list give the reason and the
  direction of the resulting error.
  \item[A7] \emph{Unpolarized.} Reflection uses the average of the two
  principal polarizations rather than tracking a polarization state.
\end{enumerate}

A6 is the assumption a reader is most likely to challenge, so it gets the space
here. Four things make it defensible at this frequency and for this quantity, and
the last of them is a bound computed on the same geometry rather than an argument
carried over from someone else's.

The first is that it has been measured, at almost this frequency and in this
kind of geometry. A street microcell campaign at 2.44, 14.8 and 58.68\,GHz, with
matched antennas and bandwidth equalized across the three bands, found the
excess loss behind a building corner growing at 3.0 to 3.5\,dB per decade of
frequency, against the 10\,dB per decade of the corner diffraction model drawn
on the same axes \cite{ericsson, mmmagic}. At one location in that campaign, at
58.68\,GHz, a manual trace of the power delay profile found no arrival above the
noise floor at the delay the path around the corner would have occupied. What
arrived first there was attributed instead to scattering from rough surfaces
near the corner, and the strongest arrival, 20\,dB above the rest of the
profile, matched a path of four specular reflections off exterior walls. That
last peak was not seen at most of the other shadowed locations, so the campaign
supports the weaker of the two readings: what reaches a shadowed receiver in a
street is not diffracted power, whether it arrives specularly or by scattering.
The mechanism is the ordinary one. A diffracted
field falls off relative to the field that produced it as $1/\sqrt{ks}$ with
distance $s$ from the edge, which at 15\,GHz and 10\,m is about $1/56$ in
amplitude, near 35\,dB down in power. Diffraction is therefore a correction to a
picture that reflection already describes. At the meter wavelengths where street
models were first built an edge is only a few wavelengths away from everything,
and the two roles are reversed.

The second is where that correction would actually lead, and it applies to one of
the two bands and not to both. Diffraction dominates only in deep shadow, at
points no direct or reflected path reaches, and the standpoints used here are
required to have open sky above them, since the standpoint filter of
Section~\ref{sec:standpoints} discards points that have none before any trace
runs. For the street cell band that closes the question, because
those sites sit from 2.5 to 6.5\,m above the head, below the roofline, and every
direction the model blocks is blocked by a facade. For the rooftop band it does
not. Those sites sit 13.5 to 43.5\,m above the head, above a typical European
roofline, and open sky above a standpoint is the precondition for a path over a
roof rather than a defense against one. Measurements in Manhattan street canyons
with a rooftop site some 20\,m above the surrounding buildings found the
over-the-top mechanism better than street level scattering by 7\,dB in fit error
\cite{adhikari}. For the rooftop band the argument therefore rests on the next
two items and not on this one.

The third is that the error has a known sign. A diffracted path can only deliver
power the trace does not currently collect, never remove any, so leaving it out
biases $\chi$ low. That holds because A2 adds powers rather than fields. Were the
estimator coherent, a diffracted contribution could arrive out of phase and the
sign argument would fail. Every value reported here is a lower bound with respect
to this term, and the shortfall grows with how narrow and how enclosed the street
is. A low bias is the right direction for an omission in an exposure study to
run, and it is stated here rather than assumed away.

% LAW CHANGE NOTE, see ../docs/LAW_CHANGE.md. The diffraction bound below is
% integrated against the illumination measure, so the 0.15 dB rooftop and
% 0.42 dB street uplifts, and the 1.14 % and 3.70 % mask disagreements, are
% old-law numbers and need recomputing. The 0.06 dB isotropic figure stands.
% The mask itself, the sign argument and the frequency argument are geometry
% and hold whatever shape Q has. The band-by-band argument two paragraphs up,
% which turns on street sites sitting below the roofline and rooftop sites
% above it, has to be made again: the new law puts sources on the facade tip
% itself, which is neither above nor below the roofline but on it.
The fourth is the size of it, measured here. A dense mask of 720 azimuths by 600
elevation bands is cast at each standpoint on the same meshes and the same
standpoint coordinates as the reported runs. Every shadowed direction is given a
Fresnel-Kirchhoff parameter $\nu = \theta\sqrt{2d/\lambda}$ and its knife edge
loss \cite{p526} is integrated against the same illumination measure the
estimator uses. The
mask first has to earn trust, and it does. It was run at the 20 lowest sky
fraction standpoints of each of three squares, 60 in all: the two most enclosed
of the eleven, Times Square and the Grand-Place, and the reference square. That
is the two most enclosed plus the reference and not the three most enclosed,
since Tokyo and Brussels both sit below the reference square on median sky
fraction. Integrated over the \emph{visible} directions instead of the shadowed
ones, its disagreement with the tracer's direct term has a median over those 60
standpoints of 0.19\% isotropic and 1.14\% rooftop, and a median of 3.70\% over
the 56 at which the street model's direct term is nonzero. The worst
disagreements over the same sets are 1.29\%, 6.44\% and 23.05\%, all of them at
single standpoints rather than at a whole square. The check uses an independent
code path with an independent quadrature. Over the same standpoints the
resulting uplift on $\chi$ has a median of 0.06\,dB isotropic, 0.15\,dB rooftop
and 0.42\,dB street. Every choice
in that calculation errs upward, so it is a bound and not an estimate: a single
absorbing half plane, a distant source, no re-blocking, a knife edge rather than
a lossy wedge, and the diffracted term multiplied by each standpoint's own
multipath gain so that paths which diffract and then reflect are covered. Against
a within square spread of 5.8 to 12.9\,dB under the rooftop model at those three
squares, 0.15\,dB changes nothing. Repeating
the same calculation at 2\,GHz raises those medians to 0.18, 0.53 and 1.78\,dB.
In linear terms the excess grows by a factor of 3.1 to 5.0, which reproduces the
frequency argument of the first item
on this geometry rather than borrowing it.

The street cell model is the exception and it is not a small one. That model puts
87.9\% of its measure below 5\textdegree{} of elevation, its worst standpoint
takes an uplift of 12.64\,dB, and at 4 of the 60 standpoints its entire elevation
support is occluded, so the direct term is exactly zero and every watt in the
answer has arrived by reflection. Those standpoints also carry the smallest
absolute $\chi$ in the study. For that model in that geometry the omitted term is
not a correction, and no argument is offered here that it is.

Two limits apply to the part of this argument that is borrowed from the
literature. One body of work fits a diffraction shaped model around a corner at
28\,GHz better than a scattering one, and read alone it points the other way.
That reading does not survive the fitted constant: the corner loss that fits is
about 2\,dB, against a theoretical edge coefficient near 40\,dB down at that
frequency \cite{duchizhik}. The functional form fits and the mechanism it is
named after is not what supplies the power. All of the published decompositions
of shadowed power into diffracted and reflected parts were measured in street
canyons, which act as waveguides, rather than in open squares. A square presents fewer facades per unit solid angle, so the
reflected budget that stands in for diffraction here is thinner than it is where
the substitution was measured. No open square decomposition exists at any
frequency to check that against. This is also where the largest untested weight
sits, since ablating diffuse scattering rather than diffraction is what moves
published prediction error the most at these frequencies \cite{vitucci}.

A7 needs a word too, because the usual defense of it is not available at this
frequency. That defense says the multiply reflected tail is depolarized, so
averaging the two polarizations is exact for it. Measured cross polarization
ratios say otherwise: pooled over many campaigns a path with no excess loss sits
near 28\,dB, falling by about half a decibel for every decibel of excess loss
\cite{karttunen}. An even split would need some 56\,dB of excess loss, and a
path that far down carries nothing. The tail that matters is not depolarized.

The cost of the average can be bounded instead, and the bound has two sides. A
vertical facade struck by a near horizontal ray presents a nearly horizontal
plane of incidence, so a vertically polarized network field meets it in the
transverse electric case, where the unpolarized mean is never wrong by more than
3.0\,dB and errs low. The transverse magnetic case has a deep minimum near
$66^\circ$ of incidence for concrete, where the mean can err high by tens of
decibels on one bounce. That geometry arises on the ground rather than on a
facade, and only for arrivals within $5^\circ$ of the $23.6^\circ$ elevation at
which a ground bounce meets that incidence, a
band holding 7.8\% of the rooftop model's measure and 0.4\% of the street
model's. So the approximation is bounded and biased low on facades, and
unbounded but rare on the ground.

The remaining sections describe the scene, the illumination model, the ray
estimator, the body, and the checks that constrain them.

\section{Configuration}
\label{sec:config}

The scene is a photogrammetric mesh of a real square, cut to a cylinder of
radius 250\,m about the standpoint and reaching from the ground to the roofline.
The mesh is built from aerial and street level imagery and carries no material
information of its own, so classes and electromagnetic parameters are assigned
afterwards. The 250\,m radius is not a free choice and Section~\ref{sec:valid}
reports what sets it. That radius should not be confused with the 250\,m that
appears in \cref{tab:models} as the far edge of the rooftop range band. One says how
much geometry is kept around the standpoint, the other how far away a base
station is allowed to be. They are fixed by different arguments and happen to
agree.

\subsection{Surfaces and materials}
\label{sec:materials}

Every triangle is assigned one of a small set of classes. In the baseline run
the class follows from geometry alone, splitting near horizontal faces at ground
level from near vertical faces above it, with vegetation separated by its
surface roughness at the mesh scale. Each class carries a complex relative
permittivity from the ITU building materials recommendation \cite{itu2040} at
15\,GHz and a root mean square roughness $\sigma_h$ from published surface
metrology on the corresponding material.

An optional path replaces the geometric class with one read from street level
photographs. Panoramas are segmented, registered to the mesh by matching the
skyline they see against the skyline the mesh casts, and their labels are
projected onto the triangles each panorama can see. Two models are used and they
answer different questions. A semantic segmentation network trained on street
scenes names the object, which separates a building from a tree from a road
surface but has one class for every facade regardless of what it is made of. A
promptable segmentation model is then asked for the material within those facade
regions, which is the axis the first model cannot resolve. Prompting a
segmentation model for a concept it was never trained to name as a class is the
idea of \cite{sam}. What ran here is SAM 3, which takes the material name itself
as the prompt.\todo{add the SAM 3 reference} The registration and projection are
described in Sections~S1 and~S2 of the supplementary information.

Imagery is taken from two providers and the split matters when reading coverage.
Squares represented by a single panorama use one provider's official captures,
whereas the multi-station route at the reference square comes from a second,
community sourced provider, which is the only set carrying the material axis.
The comparison between the geometric and the image derived assignment is a
result rather than part of the method, so it is reported separately. One thing
belongs here rather than there. The cross city comparison is run with
geometric classes at every square and carries no image evidence at all. The
photographs enter at the reference square, which is where the material path is
demonstrated and where the coverage figures of Section~\ref{sec:budget} are
measured. What limits the others is registration rather than imagery, and
until that is fixed the image derived path can claim one square.

What the photographs buy is worth stating plainly, because it is easy to expect
the wrong thing of them. They are not a fine adjustment to the permittivity.
Reflection at 15\,GHz varies little across the materials ordinary European
squares are built from, so moving between brick, stone and rendered
masonry shifts $\chi$ by far less than any difference this paper reports. The
job the photographs do is to establish which family of materials the square
belongs to at all. A square faced in glass and metal would fall outside that
family, and there the geometric default would be wrong rather than
approximate, with nothing in the mesh itself to reveal it. So the image evidence
licenses the geometric assignment instead of competing with it, and it is also
what bounds the claim: the results below are stated for masonry squares because
masonry is what was observed. It is family membership, not permittivity, that
Section~\ref{sec:budget} uses to set the bounce budget.

\subsection{Standpoints}
\label{sec:standpoints}

Exposure is not evaluated at one point. Each square contributes 80 standpoints and
results are reported as a distribution over them rather than as a single value,
because the spread within one square is comparable to the spread between
squares, so a single point would hide most of the variation.

Where a person can stand is decided from the mesh and not from a map. A grid of
columns at 3\,m spacing is laid over the crop, a ray is dropped down each one,
and the column is kept when it lands on a near horizontal face at the square's
ground level with enough clearance that a head at that position is not inside a
wall or a stall. A second filter drops any candidate whose head sees less than
1\% of the sky, which removes positions buried under an arcade or inside a
passage. That filter matters again later, because it is why no standpoint sits
in the deep shadow where the diffraction term left out by A6 would be the
leading one. Kept columns are then chained by nearest neighbor, which fixes
an order for plotting and for the step statistics but has no effect on any
exposure value, since each standpoint is traced independently. No street
centerline, footway polygon or other map layer is used, so the sampling does not
inherit whatever a mapping provider decided a street is.

The grid keeps far more candidates than are traced, so the chain is subsampled
at even spacing along its length down to the 80 that are reported. Subsampling
the chain rather than taking its first 80 links keeps the reported set spread
over the whole square instead of piled up at whichever end the chain began.

This set of standpoints is distinct from the route the photographs were captured
along, and the two should not be confused. The capture route is wherever the
imagery provider drove or walked, and it determines what the panoramas see. The
standpoints are where exposure is evaluated. They overlap but are neither equal
nor derived from one another.

\section{Illumination model}
\label{sec:illum}

% ======================================================================
% LAW CHANGE NOTE, 2026-08-03. See ../docs/LAW_CHANGE.md.
%
% This whole section describes the illumination law that is being replaced.
% It models the network by four numbers per class, a height band above the
% head and a horizontal range band, uniform in azimuth.
%
% The replacement reads the network off the geometry instead. Sites sit on
% facade tips, the top edge where a wall meets the sky, so there is no mast,
% no height band and no range cap, and each azimuth carries one source
% distance taken from the skyline rather than a distribution over one.
%
% Why: the path loss weighted branch integrates dr over [r_lo, r_hi], which
% is logarithmic in the outer cap, so d_max set the answer rather than
% bounding it. Worth 7.8 dB across the defensible span of the cap, against a
% 4.93 dB between-square spread.
%
% Nothing below has been edited. The section, tab:models and every rooftop
% and street-cell number in the results are the record of runs that happened
% under the old law. The replacement's own numbers are provisional and are
% deliberately absent.
%
% The isotropic model is untouched by the change.
% ======================================================================

Where the base stations are is unknown, so the model integrates over where they
could be. Operators do not place millimeter wave sites at arbitrary heights and
ranges: a site sits either on a roof or on street furniture, and beyond some
range its contribution stops mattering. Those two facts bound the region a site
can occupy, and what this section builds from that region is the one function the
estimator needs, the share $Q(\uext)$ of the incident power that reaches the head
from each steradian of sky.

\subsection{Sites per steradian}

Take sites spread uniformly through a slab, with $\nu$ sites per unit volume, at
heights between $h_{\min}$ and $h_{\max}$ above the head and at horizontal
ranges between $r_{\min}$ and $r_{\max}$. Look outward from $\x$ in a direction
at elevation $\alpha$, along a narrow cone of solid angle $\dOmega$. The number
of sites in that cone is the number in the part of it that lies inside the slab
and inside the range band,
\begin{equation}
  \mathrm{d}N = \nu \dOmega \int_{\rho_-(\alpha)}^{\rho_+(\alpha)} \rho^2 \,
  \mathrm{d}\rho
  = \frac{\nu}{3} \left[ \rho_+^3(\alpha) - \rho_-^3(\alpha) \right] \dOmega .
  \label{eq:cone}
\end{equation}
A site at slant range $\rho$ and elevation $\alpha$ has height $\rho \sin\alpha$
and horizontal range $\rho \cos\alpha$, so each band clips the cone from one
side,
\begin{align}
  \rho_-(\alpha) &= \max\!\left( \frac{h_{\min}}{\sin\alpha},
  \frac{r_{\min}}{\cos\alpha} \right), \nonumber \\
  \rho_+(\alpha) &= \min\!\left( \frac{h_{\max}}{\sin\alpha},
  \frac{r_{\max}}{\cos\alpha} \right).
  \label{eq:caps}
\end{align}
\Cref{fig:elevation} draws the geometry behind \eqref{eq:cone}.

What the estimator needs is not a count of sites but a share of incident power.
Under A4 the two are the same function up to a constant, because every site
contributes equally, so
\begin{equation}
  Q(\uext) \;\propto\; \rho_+^3(\alpha) - \rho_-^3(\alpha)
  \label{eq:Q}
\end{equation}
wherever $\rho_+ > \rho_-$, and $Q = 0$ elsewhere. Only elevation appears,
because a uniform density of sites on level ground has nothing to distinguish one
compass bearing from another. Counting sites rather than weighting them by range
is what makes $r_{\max}$ matter: the $\rho^2$ inside \eqref{eq:cone} says
far sites are the numerous ones, so where the range band ends decides how much of
the answer they contribute. Weighting each site by its own free space spreading
$1/\rho^2$ instead replaces the cubed bracket in \eqref{eq:Q} with
$\rho_+ - \rho_-$, and that variant is carried through as a sensitivity.

\begin{figure*}[!t]
  \centering
  \includegraphics[width=\textwidth]{18_elevation_geometry}
  \caption{A site at horizontal range $r$ and height $h$ above the head is seen
  at elevation $\alpha$ with $\tan\alpha = h/r$. Near sites appear high in the
  sky and far sites appear near the horizon, so a band of heights and a band of
  ranges together fix which elevations can hold sites at all.}
  \label{fig:elevation}
\end{figure*}

Two consequences follow directly and both matter. First, the support is not a
free parameter. A direction can hold a site only when $\rho_-(\alpha) \le
\rho_+(\alpha)$, which gives
\begin{equation}
  \arctan \frac{h_{\min}}{r_{\max}} \;\le\; \alpha \;\le\;
  \arctan \frac{h_{\max}}{r_{\min}} ,
  \label{eq:support}
\end{equation}
so writing down a height band and a range band already fixes the elevation
limits. Second, the near horizon carries most of the measure, and the near
horizon is also where the buildings act: a distant site sits low in the sky, so
its path grazes along the street and a single facade can intercept it. A site
overhead has no facade in the way. This is why the shape of the square changes
the answer at all, and it is also why the crop radius of
Section~\ref{sec:config} has to be
justified rather than chosen. \Cref{fig:box} shows the region in the height and
range plane.

\begin{figure*}[!t]
  \centering
  \includegraphics[width=\textwidth]{19_deployment_box}
  \caption{The region a site may occupy, in height above the head against
  horizontal range. Lines of constant elevation are straight through the origin,
  so the corners of the box set the elevation limits of \eqref{eq:support}.}
  \label{fig:box}
\end{figure*}

\subsection{Normalization}

$Q$ is scaled so that it integrates to one over the sphere,
\begin{equation}
  \oint Q(\uext) \dOmega = 1 ,
\end{equation}
which is what makes $\chi$ dimensionless and equal to 1 in free space. The
integral is evaluated in elevation rather than on the ray grid, because
$Q$ concentrates within a few degrees of the horizon and a few hundred
grid cells cannot resolve that. The two elevations at which \eqref{eq:caps}
changes branch are inserted into the quadrature abscissae so that no interval
straddles a kink.

The normalization hides the absolute level on purpose, which makes it easy to
forget to ask what $\chi$ multiplies. The same site density answers that. With
$\nu$ sites per unit volume of the slab, each radiating $P$, free space gives
\begin{equation}
  S_0 = \frac{P \nu}{4\pi} \oint \left( \rho_+ - \rho_- \right) \dOmega ,
  \label{eq:absolute}
\end{equation}
which is the cone integral of \eqref{eq:cone} again, now weighted by free space
spreading instead of counting sites. The $\rho_+ - \rho_-$ here is the same
object as the spreading variant of A4, used deliberately: an absolute level is a
question about how power adds, not about how a network shares it out. Over a
plausible range, 25 to 100 sites per
square kilometer radiating 55 to 75\,dBm, this puts $S_0$ between 0.0075 and
3.0\,W/m$^2$ for the rooftop band and between 0.010 and 4.1\,W/m$^2$ for the
street level band, against a general public reference level of 10\,W/m$^2$ at
this frequency \cite{icnirp}. That is nearly three orders of magnitude, set
entirely by numbers that are not public. The spread is the reason this paper reports a
ratio: the part that depends on the square is measurable and transfers between
cities, and the part that depends on the network is neither.

\subsection{Models used}

Three illumination models are carried through every run. Their parameters are in
\cref{tab:models}. The isotropic model is a control rather than a deployment: it
weights all directions equally, so under it the zero bounce part of $\chi$ is
exactly the fraction of the sphere that is open sky, which makes it a check on
the estimator as well as a reference. The full $\chi$ sits above that fraction,
by a median of 0.99\,dB over the 880 published standpoints, because the
reflected term is added to it and is positive. The rooftop model covers macro
sites on buildings and the street
model covers small cells on lamp posts and facades. Their height bands come from
the two ways operators acquire sites, a roof lease or a street furniture lease,
and published deployment heights are bimodal between them rather than spread
across the gap. Because no cellular deployment above 6\,GHz exists in the
European antenna registers, the band edges are not calibrated against measured
millimeter wave siting, and the sensitivity of $\chi$ to them is reported rather
than assumed.

% LAW CHANGE NOTE: tab:models prints the eight endpoint numbers of the old
% law. It is a model definition, so it has to be rewritten rather than
% recomputed. The isotropic row survives. See ../docs/LAW_CHANGE.md.
\begin{table}[!t]
  \caption{Illumination models}
  \label{tab:models}
  \centering
  \begin{tabular}{lccc}
    \toprule
    Model & Height above head (m) & Range (m) & Elevation (deg) \\
    \midrule
    Isotropic & -- & -- & $-90$ to $90$ \\
    Rooftop macro & 13.5 to 43.5 & 25 to 250 & 3.09 to 60.11 \\
    Street cell & 2.5 to 6.5 & 10 to 150 & 0.95 to 33.02 \\
    \bottomrule
  \end{tabular}
\end{table}

\subsection{Beam pointing}

A base station at 15\,GHz is an array, not a point radiator, so a reader will
ask what beamforming does to all of this. It depends only on whether the beam
follows the person, and the case that sounds hardest is the one that changes
nothing at all.

Suppose every site puts its peak on the pedestrian. Under A1 each source is a
plane wave from one direction, so aiming at the pedestrian means peak gain along
$-\uext$ whatever the ray did on its way out of the square. The gain toward the
head is then the same constant for every site, so it multiplies \eqref{eq:Q} and
the normalization divides it straight back out. $Q$ is not approximately
unchanged. It is the same function, and the estimator of
Section~\ref{sec:estimator} reads $Q$ only when depositing an escaped ray and
never inside a random draw, so the traced $\chi$ comes out bit identical.

For one policy this is exact and needs nothing else. Under full digital maximum
ratio transmission the site transmits along its own channel to the head, and a
matched filter delivers the total power in that channel whatever the channel is.
An array of $MN$ elements multiplies that total by $MN$ in the square and by the
same $MN$ over open ground, so the array factor cancels between numerator and
denominator with no assumption about how the paths are arranged.

The other matched policies only look invariant, and the difference matters
because they are the ones real networks use. A site that steers at the
pedestrian's coordinates aims at the direct path. A reflected path leaves that
site toward its own last surface, which can be a hundred meters from the head,
so it departs tens of degrees off the beam axis and the array puts it 20 to
25\,dB down. Geometric steering therefore suppresses the multipath this
paper is measuring and moves $\chi$ toward its no bounce value. Codebook beam
selection, which is what 5G beam management does, sits in between. Treating each
site as a single far point makes that offset vanish and both policies then read
as invariant, so the invariance there is an artifact of the far source picture
rather than a property of the physics. The direction of the error is at least
known: the numerator's gain is the peak times a sidelobe factor no greater than
one, whereas the denominator's is the peak itself, so steered $\chi$ can only
fall below the value reported here.

One thing never cancels under any policy. The element pattern multiplies each
path before the aperture combines anything, so it survives as a reweighting of
$Q$. For a standard element that is a mild elevation taper, 0.26\,dB at the
rooftop model's median elevation of $9.5^\circ$ and 10.2\,dB at the
$60^\circ$ top of the band.

What does move $\chi$ is radiation that is not aimed at the person. A broadcast
beam carries a fixed electrical downtilt, and a beam serving another user points
somewhere else entirely. In both cases the gain becomes a function of elevation
and no longer factors out of the normalization, so it reweights
\eqref{eq:Q} instead of cancelling. Since the elevation weighting is what
separates one square from another, a fixed pattern can reorder them, and that
variant is carried alongside the three models above.

\section{Adjoint ray estimator}

Reciprocity lets the trace run backwards. A ray leaving $\x$ in direction
$\uloc$, reflecting off surfaces and finally escaping in direction $\uext$ with
a surviving fraction $w$ of its launch power, describes the same physical path
as a wave arriving from a source at $\uext$ and reaching $\x$ from $\uloc$,
attenuated by the same $w$.

That reversal is what makes \eqref{eq:chidef} computable. A ray launched into
$\uloc$ and escaping into $\uext$ collects from the sources that lie in $\uext$,
which carry a share $Q(\uext)$ of the incident power, and it delivers a fraction
$w$ of it. Averaging over the rays launched into $\uloc$ gives the power arriving
from that direction, in units of $S_0$, and integrating over all launch
directions gives back the ratio itself,
\begin{equation}
  \Phi(\uloc) = \big\langle\, w \, Q(\uext) \,\big\rangle_{\uloc} , \qquad
  \chi = \oint \Phi(\uloc) \dOmega .
  \label{eq:chi}
\end{equation}
Nothing about the sources survives except $Q$, so one outward trace answers
\eqref{eq:chidef} for every source position at once. The free space denominator
$S_0$ never has to be computed, because $Q$ was normalized to make it 1.

\subsection{Estimator}
\label{sec:estimator}

The sphere of launch directions is divided into $M$ cells of near equal solid
angle using a Fibonacci grid. $N$ rays are launched in uniformly random
directions and each is assigned to the cell $c(j)$ nearest its launch direction,
giving $n_c$ rays in cell $c$. Escaping rays deposit into their launch cell, and
the estimate is the cell mean carried through the quadrature,
\begin{equation}
  \hat{\chi} = \frac{4\pi}{M} \sum_{c=1}^{M} \frac{1}{n_c}
  \sum_{j\,:\,c(j)=c} w_j \, Q(\uext_j) .
  \label{eq:estimator}
\end{equation}
Stratifying by launch cell rather than averaging over all rays at once matters
because $Q$ varies by orders of magnitude across the sphere, so an unstratified
average would be dominated by the few rays that happen to escape near the
horizon. In free space every ray escapes unobstructed with $w = 1$ and
$\uext_j = \uloc_j$, so \eqref{eq:estimator} reduces to a quadrature of $Q$ over
the sphere and returns 1, as it must.

Every ray either escapes without touching anything or reflects at least once, so
the sum in \eqref{eq:estimator} splits in two with nothing left over,
\begin{equation}
  \chi = \chi_{\mathrm{dir}} + \chi_{\mathrm{ref}} ,
  \label{eq:split}
\end{equation}
where the first term collects the untouched rays and the second collects the
rest. This separates the only two things a square can do. It hides part of the
sky, which removes direct power, and it reflects, which puts power back. The
first term is the direct arrival reduced by exactly the sky the buildings cover,
and under the isotropic model it equals the sky fraction, meaning the share of
directions from the head that reach open sky without meeting anything. That
equality is one of the identities checked in Section~\ref{sec:valid}. The second term is everything the
surroundings return. The split is a partition of the same rays, so it costs
nothing, and both parts are reported because a $\chi$ near 1 can mean an open
square or a closed one that reflects back what it blocked.

The estimator is stochastic, so every reported value carries a Monte Carlo
standard error, obtained by retracing the same standpoints over disjoint seed
streams rather than estimated from within a single run. Seeds and ray counts are
in Section~S3 of the supplementary information.

How large that error is depends strongly on which illumination model is being
evaluated, and the reason is worth stating because it decides how many rays a
result needs. The variance is set by the share of the ray budget that lands
inside the support of $Q$. The isotropic model has support everywhere, so every
ray contributes and the per standpoint error is a few thousandths of a decibel.
The street model puts its entire mass into the lowest tens of degrees of
elevation, so most rays escape into directions where $Q$ is zero and contribute
nothing, and the per standpoint error is some thirty times larger. Stratifying
the launch directions does not help with this, because the launch direction is
not the escape direction and only the latter is weighted. Reporting a median over
many standpoints does help, by roughly the square root of their number, which is
why a distribution over standpoints is the right unit for a claim and a single
standpoint often is not.

\subsection{Bounce loop}

\Cref{alg:trace} gives the loop and \cref{fig:oneray} follows a single ray
through it. Two steps deserve a word.

Roughness splits the reflection rather than blurring it. At each interaction the
fraction of reflected power that survives as a mirror reflection is the Rayleigh
factor
\begin{equation}
  \kappa = \exp\!\left[ -\left( \frac{4\pi \sigma_h \cos\theta_i}{\lambda}
  \right)^{\!2} \right] ,
  \label{eq:rayleigh}
\end{equation}
and the outgoing direction is drawn as a mirror reflection with probability
$\kappa$ and from a cosine weighted lobe about the normal otherwise
\cite{beckmann}. Choosing one lobe with a probability equal to its share of the
energy needs no compensating weight, because the two cancel. Note the
$\cos\theta_i$ in \eqref{eq:rayleigh}: a surface that scatters at normal
incidence behaves as a mirror at grazing incidence, and this is the only place
where that angle dependence enters.

Rays still traveling after $L$ interactions are dropped, and the power they
were carrying is reported as a share of the total, so the truncation is a
measured quantity rather than an assumption. Section~\ref{sec:budget} gives that
share and what follows from it.

\subsection{Bounce budget}
\label{sec:budget}

$L = 3$, and the reason is where the material evidence stops rather than where
the number stops moving. Two things have to be said straight away, because
without them this reads as a contradiction of
Section~\ref{sec:materials} and as a number chosen for cost.

The evidence in question is not a permittivity value.
Section~\ref{sec:materials} settled that those are interchangeable across
masonry. It is evidence of which family a surface belongs to, masonry against
glass or metal, which is the one material distinction that does move $\chi$ and
the one the mesh cannot supply on its own. So the budget is set by where
observed family membership stops, not by where a number stops being known to
three decimals.

The evidence does not reach indefinitely either, and the measurements below say
where it runs out rather than leaving that to a reader to wonder about. A
convergence threshold would be the usual argument in place of all this, and it is
not weaker so much as insufficient by itself, since it shows the discarded tail
is small without showing the kept part is right. Both arguments are used here,
and the convergence one comes last.

Two different questions hide inside the phrase ``observed surface'', and they
have to be separated or the numbers below look like they disagree. The first is
geometric: can the observation point see this surface at all? The second is
evidential: of the surfaces it can see, which did a camera actually photograph?
The first is a property of the formulation and holds in any square. The second
is a property of this particular set of panoramas.

Take the geometric question first. A photograph taken from the observation point
sees exactly the surfaces a ray can reach on its first interaction, so the first
hit is covered by construction. For a path that returns to where it started the
cover reaches the second interaction as well, since both ends of a closed loop
are seen from the same viewpoint. This is the co-located picture of
Section~\ref{sec:ratio} in its strictest form, and it is not a hypothetical.
Traced at 42 standpoints across four photogrammetric meshes, at the Grand-Place,
the reference square, Milan and Tokyo,
with the visible set built by first hit probing rather than assumed, the median
share of returning power whose whole chain lies on visible surfaces is 1.000000
at the first order and 0.999972 at the second. It falls to 0.980 at the third,
because the cover is over the two ends of a loop and not over its middle. The
spread across the four cities is a few percent, so the first two orders are a
theorem and the third is a property of the scene.

An outward path gives up half of that, and this is what the adjoint move costs.
A ray leaving the first facade travels away from the viewpoint that saw it, so
only its first hit is assured. On the same geometry an outward path holds
0.999103 at the first interaction, 0.912004 at the second and 0.714182 at the
third, so by the second interaction the estimator is already carrying 9\% of
its power over surfaces the standpoint cannot see, and by the third 29\%.

The evidential question is the stricter one, because a camera has to have stood
somewhere and looked. Where one did, the two questions coincide: the share of
incident power landing on a triangle some panorama photographed is 0.978 to
0.999 at the first interaction, across seven of the eight registered stations at
the reference square, which confirms the geometric argument to about four parts
in a thousand at the median and twenty two at the worst of the seven. The eighth
returns 0.362, and since it stands eight meters from
one scoring 0.999 that is a registration failure rather than a coverage one. It
had passed the skyline residual gate, so a pose can clear the residual and still
fail on coverage. First interaction coverage is therefore worth reading beside
that residual, it costs one trace, and it is available at any square with a
registered panorama. No claim is made here about which step of the registration
produced the bad pose: the probe standpoints behind these eight figures are
provider poses displaced 3.9 to 5.5\,m from their registered cameras, and the
displacement does not order the coverage, since the one displaced furthest
scores 0.978. The station is excluded from the coverage figures above, and no
station of the nine in the 250\,m binding falls below 0.996. Away from a camera the two
questions come apart, and the reach turns out to be a radius. Within 20\,m of a
registered panorama the second interaction still lands on photographed material
93 to 96\% of the time and the third about 89\%, whereas beyond 40\,m no
depth reaches 0.22. Pooled over the published walk, whose standpoints are
sampled within a 90\,m radius and whose cameras span 56\,m end to end, first
interaction coverage is 0.479. So the evidential claim is about a twenty meter
neighborhood of the instrument rather than about a square, and every result
that uses image derived material carries that number beside it. These
percentages and the geometric ones above come from different standpoint sets at
different crop radii, so they answer the same question in different places and
should not be subtracted from one another.

Three is then the smallest budget that lets both questions be asked. One
interaction would test nothing beyond the assured hit, two would show whether
the walk covers the square, and three shows what the first clearly unobserved
interaction is worth.

It matters which of the two legs is doing the licensing, and where. The
evidential leg is measured at the one square that has registered panoramas along
a walk, so on its own it licenses $L = 3$ only there. At the other ten squares
the budget rests on the power leg and the convergence measurement below, which
are computed at every square and do not need a camera. The two legs are not
independent of each other in the way that would be convenient: they agree here,
but agreement at one square is not proof that the evidential reach is the same
elsewhere. What can be said across squares is that the geometric visibility
above was measured on four meshes and varied by a few
percent, and that power-weighted photographic coverage at a station position,
pooled over the stations of a square, runs from 0.941 to 0.999 at the first
interaction over a second set of four squares holding 35 stations between them.
Pooling is what makes that range narrow: taken one station at a time the same
35 run from 0.696 to 0.9998. The reported budget is therefore set by evidence at
one square and carried to the others by power, with the geometric part checked
in four.

What makes three sufficient rather than merely defensible is where the power is.
The shares that follow were measured at the reference square at the eight
registered station positions on a 130\,m crop, which is not the 250\,m crop the
reported runs use, so the 250\,m value is given beside each one. Of all launched
power, 21.0\% never touches a surface and escapes straight
to the sky when pooled over those eight camera positions, 23.6\% when pooled
over the 40 walk standpoints instead and 23.3\% at the wider crop, which is the
direct term of \eqref{eq:split}. The power reaching the first, second and third surfaces is
79\%, 8.8\% and 1.2\% of the launched total, and 76.7\%, 8.7\% and 1.0\% at the
wider crop. Those shares are not parts of one whole and are not to be added,
because power that reaches the third surface was already counted at the first
and the second. Each is what is left of the same launched power at a greater
depth. Interactions past the third carry 0.23\% of launched power, 0.17\% at the
wider crop. Those depths can be justified
neither by evidence nor by the power they carry. $L = 3$ costs a median of
0.002\,dB against an eight interaction reference, at fixed seeds so that only
the budget moves. The worst standpoint it moves depends on which square is
asked. Over the 40 standpoints of the published walk at the reference square and
a 130\,m crop the worst is 0.063\,dB; over four squares at the 250\,m crop the
median moves by at most 0.021\,dB and the worst single standpoint by
0.384\,dB, at the Grand-Place under the street model. The bound that holds
everywhere it was measured is the second one. Not one standpoint of either set
moves half a decibel under any of the three illumination models, which is what
the budget argument rests on.

One consequence has to be stated plainly. A fixed budget is a hard cut, so the
estimator is truncated and therefore biased low, since truncation deletes power
and never adds it. The share of escaping power still in flight when the budget
runs out is 0.0037 at the median of those 40 standpoints and 0.0201 at the
worst. Over the 880 standpoints of the reported runs it is 0.0050 at the median
and 0.0795 at the worst, again at the Grand-Place, and every
value reported here is a lower bound in that sense. This is not the same number
as the 0.23\% above, because the denominators differ: one is a share of all
launched power and the other a share of the power that escapes. Russian roulette is the
standard way to end paths without introducing this bias \cite{arvokirk}, but it
cannot repair a hard cap, and at a budget of three it would fire exactly once,
one step before the cap, killing rays at a 0.05 survival floor and inflating the
few survivors twentyfold. An eight seed test over four standpoints found the
relative spread agreeing to two significant figures on the isotropic and rooftop
models and to three on the street model with roulette on and off, and wall times
of 101 and 100\,s. It is switched off because it changes neither variance nor
work at this budget. The only approximation left in the trace itself is a single one sided
truncation, whose size is quoted above. The assumptions of
Section~\ref{sec:ratio} sit underneath that and are not part of the claim.

\begin{algorithm}[!t]
\caption{Adjoint trace at one standpoint}
\label{alg:trace}
\begin{algorithmic}[1]
\State $\hat{\Phi} \gets 0$
\For{$j = 1$ to $N$}
  \State draw $\uloc_j$ uniformly on the sphere, $c \gets$ nearest grid cell
  \State $w \gets 1$, $\uext \gets \uloc_j$
  \For{$b = 0$ to $L$}
    \State intersect the scene along $\uext$
    \If{no hit}
      \State $\hat{\Phi}[c] \gets \hat{\Phi}[c] + w\,Q(\uext)$
      \State \textbf{break}
    \EndIf
    \State $w \gets w \, |\Gamma(\theta_i, \varepsilon_r)|^2$
    \State with probability $\kappa$ set $\uext$ to the mirror direction,
    \Statex \hspace{\algorithmicindent} otherwise draw it cosine weighted about $\nhat$
  \EndFor
  \State \Comment{the loop ends at $b = L$ and whatever is still
  \Statex \hspace{\algorithmicindent} in flight is discarded}
\EndFor
\State \Return $\hat{\chi}$ from \eqref{eq:estimator}
\end{algorithmic}
\end{algorithm}

The cost is $O(N L \log T)$ for $T$ scene triangles, the logarithm coming from
the bounding volume hierarchy. Memory is $O(M)$ per standpoint, since no path is
stored: the trace keeps angular histograms and scalar sums only.

\begin{figure*}[!t]
  \centering
  \includegraphics[width=\textwidth]{21_one_ray}
  \caption{One ray leaving the head, reflecting twice off facades and escaping
  to the sky. Line thickness is the surviving power $w$. The ray deposits
  $w\,Q(\uext)$ into the angular cell it was launched into, not the one it
  escaped through.}
  \label{fig:oneray}
\end{figure*}

\subsection{Quantities that come for free}

Two further quantities fall out of the same trace at no extra cost, on top of
the split in \eqref{eq:split}. The sky fraction, already met above, is reported
on its own because it depends on the geometry alone and on no illumination
model, and it is reported alongside $\chi$ throughout. The mean excess path
length, taken over escaping rays weighted by their power, gives the delay spread
the square would impose. That one is recorded per standpoint but is not analyzed
in this paper, and it is mentioned here only because it costs nothing to keep.

\section{Body coupling}

The trace gives the power arriving at $\x$ from each direction. Turning that
into absorbed power on a body is a separate calculation and is done with an
established phantom and tissue model rather than reimplemented here
\cite{aegis, itis}. For a facet with outward normal $\nhat$ under a plane wave
traveling along $\khat$,
\begin{equation}
  S_{ab} = S_{\mathrm{inc}} \, T_0 \, \max\!\left[ \nhat \cdot (-\khat),\, 0
  \right] ,
  \label{eq:sab}
\end{equation}
where $T_0$ is the power transmitted into skin at normal incidence at 15\,GHz
and the clamp removes facets that face away. The arriving distribution
$\Phi(\uloc)$ is passed in as a set of plane waves, one per angular cell, with
$\khat = -\uloc$, as reciprocity requires, and each carrying
$S_{\mathrm{inc}} = S_0 \, \Phi(\uloc) \, \Delta\Omega$ over a cell of solid
angle $\Delta\Omega$. Summing
\eqref{eq:sab} over cells gives the absorbed power density map on the body, and
integrating that map gives the absorbed power and the whole body specific
absorption rate.

Two things are left out of this step, and the first matters more than its size
suggests. The phantom stands alone. A bystander is not a small perturbation at
15\,GHz: a person is roughly thirty wavelengths across and optically thick, so
another body is opaque geometry rather than something a wave works around. What
makes a crowd different from a building is that it stands at the head's own
height. A standing adult reaches about 1.73\,m against an observation height of
1.5\,m, so a bystander adds only some 0.23\,m above the head, and a ray leaving
at elevation $\alpha$ clears the crowd after a horizontal distance
$0.23\,\mathrm{m} / \tan\alpha$. That is 13\,m at $1^\circ$, 2.6\,m at
$5^\circ$ and 0.4\,m at $30^\circ$. A crowd is therefore opaque along the
horizon and transparent above roughly $10^\circ$, and no crowd density changes
that.

The obvious next step is to predict each model's sensitivity from how much of
its measure sits near the horizon, and it is instructive that this fails. The
street level model sends 88\% of its measure below $5^\circ$ and the rooftop
model 9\%, a ratio of nine. The shares that actually arrive at the head are
0.212 and 0.185, a ratio of 1.15. The facades have already redistributed the power
before any crowd sees it, so a dense European square is its own near horizon
blocker and a crowd is a second edit to a distribution the buildings have
already set. Measured at 12 standpoints against paired empty square
baselines, 2.15 people per square meter of walkable ground cost a median
1.40\,dB under the street
model and 1.07\,dB under the rooftop model, and an ordinary busy 0.3 per square
meter costs 0.5 and 0.3\,dB. The isotropic control loses only 0.07\,dB, and
the reason is that weighting every direction equally makes redistribution free
and leaves only absorption to pay for. This mechanism was measured directly:
repeating the run with the bodies replaced by index matched absorbers, so that
the geometry and the blockage are identical and only the reflected half is
removed, takes the isotropic loss from 0.07 to 0.71\,dB while the rooftop and
street columns move only from 1.07 to 1.33 and from 1.40 to 1.61\,dB. That
control ran at 8 standpoints and one crowd realization against the 12 and two of
the crowd arm, sharing 2 standpoints with it, so the pairs above compare two
populations and are a contrast in mechanism rather than a difference of one
number. A crowd
is nearly invisible to the isotropic control because it reflects, not because it
fails to block.

The second omission is that the body is treated as convex, so one limb never
shadows another, which matters most for arrivals that graze along it.

\section{Validation}
\label{sec:valid}

The estimator is checked in three ways, ordered from those that need no
reference value to those that need a reference calculation.

\subsection{Identities}

Some results follow from the construction and must come out exactly. Under the
isotropic model the zero bounce term of $\chi$ must equal the sky fraction,
because weighting all directions equally makes that term the open solid angle
and nothing else, and the two are computed by different code paths. Over the 880
published standpoints the two agree to a median of 0.0004\,dB and never differ
by more than 0.0651\,dB. The identity holds for the zero bounce term and not for
$\chi$, which also carries the reflected term and therefore exceeds the sky
fraction by a median of 0.99\,dB and by as much as 2.94\,dB. In an empty scene
every model must return $\chi = 1$. Reversing every path must leave the
transported power unchanged. Total reflected power must never exceed incident
power at any interaction. These are property tests over randomized inputs rather
than single cases.

\subsection{Closed forms}

Two geometries have analytic answers. A single infinite ground plane below the
standpoint gives a $\chi$ that can be written down from \eqref{eq:cone} and the
Fresnel coefficient alone. A Lambertian spherical cavity gives a second, with a
different balance between the direct and the multiply reflected contribution.
Both are reproduced to the tolerance set by the ray count.

\subsection{Convergence}

Two parameters are measured here, and a third is corroborated
rather than chosen. The bounce budget $L$ is fixed by evidence in
Section~\ref{sec:budget}, and the trace out to eight interactions reported
there is what corroborates it. That comparison needs one remark, because the obvious way to read it is
wrong. Two runs that differ only in $L$ and share a seed trace the same rays and
agree exactly up to the shallower depth, so their difference is paired and is
resolved far better than either value is on its own. Reading it against the single run
error of Section~\ref{sec:estimator} would make a real and well measured
difference look like noise. The ray count $N$ is raised until the spread across
independent seeds falls below the differences being interpreted. The crop radius is raised until
$\chi$ stops moving, which is what fixes the 250\,m figure quoted in
Section~\ref{sec:config}: at smaller radii the mesh is missing buildings that
would have
blocked near horizon paths, and the near horizon is where the illumination
models put most of their measure. The converged radius is a property of the
illumination model and not of the square, so it is reported per model.

% LAW CHANGE NOTE, see ../docs/LAW_CHANGE.md. Three things in this paragraph
% and the one below go with the old law.
%  1. The converged crop radius per model, and the claim that 250 m follows,
%     rest on the street model's low elevation measure. Remeasure.
%  2. The whole rho_+ paragraph sweeps a parameter the new law does not have.
%     The 7.4 dB and the rho_+^-1.2 fit stay as the diagnosis, not as a
%     sensitivity band. The point that no absolute chi can be quoted without
%     its radius is exactly what the change is meant to remove.
%  3. "every band law is uniform in azimuth, so one trace reduces to an
%     elevation histogram" does NOT carry over. The facade tip law reads a
%     skyline, so it is not azimuth uniform and the reduction is lost. This is
%     machinery, not just a number.
One radius converges and one does not, and they are easy to confuse. The crop
radius above is the extent of the mesh. The outer radius $\rho_+$ of the
deployment box in Section~\ref{sec:illum} is a different number, and sweeping it
from 100 to 400\,m never reaches a plateau: $\chi$ falls about as
$\rho_+^{-1.2}$, worth 7.4\,dB at the median square for the rooftop model. This
is not a defect of the estimator. It is what a uniform site density over an
unbounded plane does, since the lower edge of the elevation support keeps
sliding toward the horizon. It means no absolute $\chi$ can be quoted without
the radius it was computed at. What is far steadier is the contrast between
squares, which moves 0.51\,dB against that 7.4\,dB, and the results lead with
that contrast. The sweep costs no extra tracing, for the same reason the
adjoint move works: $\chi$ is linear in the illumination density and every band
law is uniform in azimuth, so one trace reduces to an elevation histogram and
any choice of box is a dot product against it. Every radius is therefore scored
on identical rays.

Identities and closed forms constrain the implementation, not the formulation.
A derivation error that is coded faithfully passes all of them, and one such
error was in fact found by rederiving \eqref{eq:cone} rather than by any test
failing. What guards against that class of error is a second calculation that
shares no code with the first. The diffraction bound of A6 provides one as a
by-product: its mask, its quadrature and its ray budget are independent of the
estimator. Integrated over the visible directions, its disagreement with
the estimator's direct term has a median over the 60 standpoints of
Section~\ref{sec:ratio} of 0.19\% isotropic and 1.14\% rooftop, and a median of
3.70\% over the 56 at which the street model's direct term is nonzero. Only the
direct term is covered by this check. The reflected
term is checked against Sionna RT, an independently written tracer, run on the
same meshes with the same materials. The worst of the three illumination models
agrees to 0.25\,dB at the reference square and 0.16\,dB at the Grand-Place, each
being a median over standpoints. The median across the three models is 0.24 and
0.03\,dB. The Monte Carlo floor on the comparison itself is 0.23 to 0.35\,dB in
the configuration reported here. The comparison therefore resolves no
disagreement.

The reflected check has two limits. It cannot test the illumination law because
$Q$ is common to both calculations and an error in it cancels exactly, so it
tests transport rather than weighting. The two calculations are also matched on
roughness rather than on model. This estimator forms a product of per-interaction
polarization averages, whereas the other forms an average of products. The
residual above contains that difference rather than excluding it.
